\chapter*{외부 검토용 판본 안내}
\addcontentsline{toc}{chapter}{외부 검토용 판본 안내}
\markboth{외부 검토용 판본 안내}{외부 검토용 판본 안내}

이 문서는 Open Algebra KR Volume 3의 \BookVersion, \BookRevision이다. Chapter 0--8의 전권 초고와 두 차례의 내부 통합 검수를 마친 뒤, 독립적인 수학 검토와 전문 한국어 교열·판면 검수를 받기 위해 공개한 판본이다.

\begin{keyidea}
이 판본의 공개는 외부 검토가 완료되었다는 뜻이 아니다. 저자와 AI 도구가 수행한 내부 검수와, 집필에 참여하지 않은 독립 검토자의 확인은 구분하여 기록한다.
\end{keyidea}

\section*{검토를 요청하는 항목}

\begin{enumerate}[label=(\arabic*)]
  \item \textbf{수학적 정확성}: 정의의 가정, 정리의 범위, 증명의 누락, 예제와 연습문제의 계산을 확인한다.
  \item \textbf{학습 흐름}: 선수개념의 배치, 장간 의존관계, 같은 정리가 다시 등장하는 이유가 충분히 설명되는지 확인한다.
  \item \textbf{한국어 교열}: 문장 호응, 띄어쓰기, 용어의 일관성, 영문 고유명사와 수식 주변의 문장을 확인한다.
  \item \textbf{판면과 접근성}: 글자·수식·표·색상 상자의 가독성, 페이지 나눔, 목차와 링크, 흑백 출력에서의 구별 가능성을 확인한다.
\end{enumerate}

\section*{오류를 보고하는 방법}

프로젝트 저장소의 Issues에서 검토 종류에 맞는 양식을 사용한다.
\begin{center}
  \url{https://github.com/SengangLemon/open-algebra-kr/issues}
\end{center}
보고에는 가능한 한 다음 정보를 포함한다.

\begin{itemize}
  \item 판본: \BookVersion
  \item 물리적 PDF 페이지와 장·절·정리 또는 문제 번호
  \item 문제가 되는 문장이나 수식의 짧은 인용
  \item 오류의 이유와 재현 가능한 계산 또는 참고문헌
  \item 제안하는 수정과 심각도: 치명적, 주요, 경미, 교열, 판면
\end{itemize}

검토 지침과 장별 고위험 항목은 저장소의
\texttt{notes/reviews/volume-03-external-review-guide.md}에 정리되어 있다. 접수된 지적은 공개 이슈와 수정 커밋으로 추적하며, 독립 검토와 전문 교열·디자인 검수가 끝난 뒤에만 릴리스 후보로 승격한다.
